% from ega2-1-fr.tex
\oldpage[II]{202}
$\mathcal{S}$-Module gradué quasi-cohérent $\mathcal{M}$ et tout entier $k$
tel que $0\leqslant k\leqslant d-1$, on a aussi un $h$-isomorphisme canonique
(avec les notations de (\hyperref[II.3.1.1-fr]{3.1.1}))
\begin{equation}
\label{II.8.14.13.1-fr}
  (\operatorname{Proj}(\mathcal{M}))^{(d,k)}
  \xleftarrow{\sim}\operatorname{Proj}(\mathcal{M}^{(d,k)}).
\tag{8.14.13.1}
\end{equation}

Supposons d'abord $Y=\operatorname{Spec}(A)$ affine,
$\mathcal{S}=\widetilde{S}$, $\mathcal{M}=\widetilde{M}$, où $S$ est une
$A$-algèbre graduée à degrés positifs et $M$ un $S$-module gradué. On sait
que pour tout $f\in S_e$ ($e>0$) $h$ applique $D_+(f)$ sur $D_+(f^d)$ et
correspond à l'isomorphisme canonique
$S_{(f^d)}\xrightarrow{\sim}S_{(f)}$
(\hyperref[II.2.2.2-fr]{2.2.2}). La restriction de
(\hyperref[II.8.14.13.1-fr]{8.14.13.1}) à $D_+(f^d)$ correspond alors au
di-isomorphisme canonique $M_{f^d}\xrightarrow{\sim}M_f$, restreint aux
éléments de $M_{f^d}$ de degré congru à $k$ (mod. $d$). Nous laissons au
lecteur le soin de vérifier que ces isomorphismes sont compatibles avec le
passage de $f$ à un de ses multiples homogènes $fg$, puis qu'on a une
compatibilité analogue avec le passage de $S$ à une $A'$-algèbre graduée
$S'=S\otimes_A A'$, lorsque $A'$ est une $A$-algèbre.

En particulier, on a ainsi un $h$-isomorphisme
\begin{equation}
\label{II.8.14.13.2-fr}
  (\mathcal{S}^{(d)})_{X^{(d)}}
  \xrightarrow{\sim}(\mathcal{S}_X)^{(d)}
\tag{8.14.13.2}
\end{equation}
qui respecte les structures multiplicatives des deux membres, et grâce auquel
(\hyperref[II.8.14.13.1-fr]{8.14.13.1}) devient un $h$-di-isomorphisme d'un
$(\mathcal{S}^{(d)})_{X^{(d)}}$-Module gradué sur un
$(\mathcal{S}_X)^{(d)}$-Module gradué. De même, on a un $h$-isomorphisme
\begin{equation}
\label{II.8.14.13.3-fr}
  \operatorname{Proj}_0(\mathcal{S}^{(d,k)}(n))
  \xrightarrow{\sim}\mathcal{O}_X(nd+k)
\tag{8.14.13.3}
\end{equation}
qui complète le résultat de (\hyperref[II.3.2.9-fr]{3.2.9, (ii)}).

On déduit immédiatement de l'isomorphisme
(\hyperref[II.8.14.13.1-fr]{8.14.13.1}) un isomorphisme de
$\mathcal{S}^{(d)}$-Modules gradués
\begin{equation}
\label{II.8.14.13.4-fr}
  \Gamma_*^{(d)}(\operatorname{Proj}(\mathcal{M}^{(d,k)}))
  \xrightarrow{\sim}
  \Gamma_*((\operatorname{Proj}(\mathcal{M}))^{(d,k)})
\tag{8.14.13.4}
\end{equation}
où $\Gamma_*^{(d)}$ correspond au morphisme structural
$q^{(d)}:X^{(d)}\to Y$~; il est immédiat de vérifier que l'homomorphisme
canonique $\alpha$ (\hyperref[II.8.14.5.4-fr]{8.14.5.4}) et l'homomorphisme
analogue $\alpha^{(d)}$ pour $X^{(d)}$ rendent commutatif le diagramme
\begin{equation}
\label{II.8.14.13.5-fr}
  \xymatrix{
    & \mathcal{M}^{(d,k)}\ar[dl]_{\alpha^{(d)}}\ar[dr]^{\alpha} & \\
    \Gamma_*^{(d)}(\operatorname{Proj}(\mathcal{M}^{(d,k)}))
      \ar[rr]^{\sim}
      && \Gamma_*((\operatorname{Proj}(\mathcal{M}))^{(d,k)})
  }
\tag{8.14.13.5}
\end{equation}
la vérification se faisant en supposant $Y$ affine, et calculant les
restrictions des images par $\alpha^{(d)}$ et $\alpha$ d'un même élément de
$M^{(d,k)}$ aux ouverts $D_+(f^d)$ et $D_+(f)$ (avec les mêmes notations que
ci-dessus).
\end{remark}

\begin{proposition}[8.14.14]
\label{II.8.14.14-fr}
Soient $Y$ un préschéma quasi-compact, $\mathcal{S}$ une
$\mathcal{O}_Y$-Algèbre graduée quasi-cohérente de type fini, $\mathcal{M}$
un $\mathcal{S}$-Module gradué quasi-cohérent vérifiant la condition
\textbf{(TF)}~; soit $X=\operatorname{Proj}(\mathcal{S})$. Alors
$\mathcal{S}_X$ est une $\mathcal{O}_X$-Algèbre graduée périodique
(\hyperref[II.8.14.12-fr]{8.14.12}), et il existe une

